\chapter{Patching and Diffing: Updating the DOM Efficiently}

Mounting is the ``first render'' story. But real apps are dominated by \emph{updates}: state changes, input events, list reorders, and component re-renders.

Relax's patcher has one job: given an \emph{old} VNode tree that is already mounted and a \emph{new} VNode tree produced by the latest render, mutate the real DOM so it matches the new tree while reusing as much existing DOM as is safe.

In Relax the patcher is both:

\begin{itemize}
  \item \textbf{small enough to read end-to-end} (in \texttt{src/patch-dom.ts}), and
  \item \textbf{heavily specified by tests} (especially \texttt{src/\_\_tests\_\_/patch-dom.test.ts}).
\end{itemize}

This chapter is long because patching is where most correctness and performance edge cases live.

\section{Diagram: the patch decision tree}
Before reading any helper functions, keep this diagram in mind. It summarizes the control flow inside \texttt{patchDOM}:

\begin{center}
\begin{tikzpicture}[
  node distance=8mm,
  box/.style={draw,rounded corners,fill=RelaxLight,inner sep=5pt,align=left},
  arrow/.style={-Latex,thick}
]
  \node[box] (start) {\texttt{patchDOM(oldVdom, newVdom, parentEl, hostComponent?)}};
  \node[box, below=of start] (eq) {\textbf{Decision:} \texttt{areNodesEqual(oldVdom, newVdom)}?};

  \node[box, below left=of eq, xshift=14mm] (replace) {\textbf{Replace path}\\
  	exttt{index = findIndexInParent(parentEl, old.el)}\\
  	exttt{destroyDOM(old)}\\
  	exttt{mountDOM(new, parentEl, index, hostComponent)}};

  \node[box, below right=of eq, xshift=-14mm] (patch) {\textbf{Patch path}\\
  	exttt{new.el = old.el}\\
  	exttt{switch (new.type) \{ ... \}}};

  \node[box, below=of patch] (dispatch) {\textbf{Per-type patchers}\\
  TEXT \textrightarrow{} patchText\\
  ELEMENT \textrightarrow{} patchElement\\
  FRAGMENT \textrightarrow{} patchFragmentChildren\\
  COMPONENT \textrightarrow{} patchComponent\\
  HRBR \textrightarrow{} patchHrbr};

  \draw[arrow] (start) -- (eq);
  \draw[arrow] (eq) -- node[above left]{no} (replace);
  \draw[arrow] (eq) -- node[above right]{yes} (patch);
  \draw[arrow] (patch) -- (dispatch);
\end{tikzpicture}
\end{center}

This picture also explains why Chapter~4's invariant (children are always VNodes) matters: it makes the \texttt{switch} and its helper patchers tractable.

\section{\texttt{patchDOM}: the top-level operation}
The public entry point is:

\begin{quote}
    exttt{patchDOM(oldVdom, newVdom, parentEl, hostComponent?)}
\end{quote}

It lives in \texttt{src/patch-dom.ts}.

\subsection{Contract}

\begin{itemize}
  \item \textbf{Inputs}
    \begin{itemize}
      \item \texttt{oldVdom}: the currently mounted VNode tree.
      \item \texttt{newVdom}: the VNode tree produced by the latest render.
      \item \texttt{parentEl}: the DOM parent that contains \texttt{oldVdom.el}.
      \item \texttt{hostComponent}: optional component instance used to bind events and to account for fragment offsets.
    \end{itemize}
  \item \textbf{Output}
    \begin{itemize}
      \item Returns the patched \texttt{newVdom} (with \texttt{newVdom.el} pointing at the reused DOM).
    \end{itemize}
  \item \textbf{Key invariant}
    \begin{itemize}
      \item After patching, the returned tree (the \emph{new} VDOM) must contain correct \texttt{.el} references.
    \end{itemize}
\end{itemize}

\subsection{Replace vs patch: areNodesEqual}

The first decision is whether two nodes are ``the same'' (so we can patch in place) or different (so we must replace).

Relax defines ``same'' in \texttt{src/nodes-equal.ts} via \texttt{areNodesEqual(oldNode, newNode)}.

This is not a philosophical definition of equality; it is \emph{exactly} the predicate the patcher uses to decide between:

\begin{itemize}
  \item patching in place (reuse), or
  \item destroying and remounting (replace).
\end{itemize}

The current identity rules are:

\begin{itemize}
  \item If \texttt{type} differs: not equal.
  \item \textbf{TEXT}: always equal (text nodes are patched by content).
  \item \textbf{ELEMENT}: equal when \texttt{tag} matches and \texttt{key} matches.
  \item \textbf{COMPONENT}: equal when component function identity (\texttt{tag}) matches and \texttt{key} matches.
  \item \textbf{HRBR}: equal when \texttt{mount} factory identity matches.
\end{itemize}

\paragraph{Why identity is intentionally strict}
Virtual DOM patching is always a trade-off between reuse and correctness. Relax chooses a conservative identity rule:

\begin{itemize}
  \item If identity says "not equal", Relax pays the cost of teardown + remount.
  \item The upside is that "equal" becomes a promise: the patcher can safely assume the DOM node kind matches what it needs to mutate.
\end{itemize}

In other words, \texttt{areNodesEqual} isn't about "deep equality". It's a \emph{safety gate} that defends the patcher.

\paragraph{A renderer mental model: identity answers ``can I mutate in place?''}
When I say two VNodes are ``equal'' in Relax, I don't mean their props match.
I mean something narrower and more operational:

\begin{quote}
If I keep the existing DOM node, do I still have the right \emph{kind} of DOM node to mutate into the new node?
\end{quote}

That’s why:

\begin{itemize}
  \item TEXT nodes are always patchable (you can always update \texttt{nodeValue}).
  \item ELEMENT nodes require the same \texttt{tag} and the same \texttt{key}.
  \item COMPONENT nodes require the same component function identity (\texttt{tag}) and the same \texttt{key}.
\end{itemize}

That strictness is what keeps the rest of the patcher simple.
When identity says ``yes'', every per-type patch function can assume it's operating on the right DOM shape.

If nodes are not equal, \texttt{patchDOM} does a replacement:

\begin{enumerate}
  \item Find the current DOM index of \texttt{oldVdom.el} inside \texttt{parentEl} (\texttt{findIndexInParent}).
  \item Destroy the old subtree (\texttt{destroyDOM(oldVdom)}).
  \item Mount the new subtree at the same index (\texttt{mountDOM(newVdom, parentEl, index, hostComponent)}).
\end{enumerate}

This behavior is covered early in \texttt{src/\_\_tests\_\_/patch-dom.test.ts} by:

\begin{itemize}
  \item \texttt{change the root node}
  \item \texttt{change the root node, mount in same index}
\end{itemize}

\subsection{When nodes are equal: preserve el}

If nodes are equal, patching starts by reusing the DOM reference:

\begin{lstlisting}[language=JavaScript]
newVdom.el = oldVdom.el
\end{lstlisting}

This is asserted by tests like \texttt{no change} and \texttt{sets the el in the new vdom}.

\section{Node-type specific patching}

After equality is established, patching dispatches on \texttt{newVdom.type} (see \texttt{DOM\_TYPES} in \texttt{src/h.ts}).

\subsection{Text nodes: patchText}

\texttt{patchText(oldVdom, newVdom)} compares \texttt{oldVdom.value} and \texttt{newVdom.value}. If different, it updates \texttt{Text.nodeValue}.

Test: \texttt{patch text}.

\subsection{Element nodes: patchElement}

\texttt{patchElement(oldVdom, newVdom, hostComponent)} updates four categories of props:

\begin{itemize}
  \item Attributes/properties (everything except \texttt{class}, \texttt{style}, \texttt{on})
  \item \texttt{class}
  \item \texttt{style}
  \item \texttt{on} (DOM event handlers)
\end{itemize}

Implementation detail: it destructures props like:

\begin{lstlisting}[language=JavaScript]
const { class: oldClass, style: oldStyle, on: oldEvents, ...oldAttrs } = oldVdom.props
const { class: newClass, style: newStyle, on: newEvents, ...newAttrs } = newVdom.props
\end{lstlisting}

and then applies targeted patchers:

\begin{itemize}
  \item \texttt{patchAttrs(el, oldAttrs, newAttrs)}
  \item \texttt{patchClasses(el, oldClass, newClass)}
  \item \texttt{patchStyles(el, oldStyle, newStyle)}
  \item \texttt{patchEvents(el, oldListeners, oldEvents, newEvents, hostComponent)}
\end{itemize}

\paragraph{Micro-optimization: object identity fast paths}

Notice the checks in \texttt{patchElement}:

\begin{verbatim}
if (oldAttrs !== newAttrs) { ... }
if (oldClass !== newClass) { ... }
if (oldStyle !== newStyle) { ... }
\end{verbatim}

In large lists, applications often reuse the same props objects (or stable sub-objects) for unchanged rows.
These identity checks avoid unnecessary diff work.

\subsection{Attributes: patchAttrs uses objectsDiff}

\texttt{patchAttrs} calls \texttt{objectsDiff(oldAttrs, newAttrs)} (from \texttt{src/utils/objects.ts}) to compute \texttt{added}, \texttt{removed}, \texttt{updated}.

It then:

\begin{itemize}
  \item removes: \texttt{removeAttribute(el, name)}
  \item adds/updates: \texttt{setAttribute(el, name, value)}
\end{itemize}

Important: \texttt{setAttribute} and \texttt{removeAttribute} are the same helpers discussed in Chapter 6 (\texttt{src/attributes.ts}), including the property-vs-attribute rules.

\subsection{Classes: patchClasses uses arraysDiff}

\texttt{patchClasses} normalizes both old and new values into an array of class tokens using \texttt{toClassList}.

Then it runs \texttt{arraysDiff(oldClasses, newClasses)} to compute \texttt{added} and \texttt{removed}, and performs:

\begin{itemize}
  \item \texttt{el.classList.remove(...removed)}
  \item \texttt{el.classList.add(...added)}
\end{itemize}

The behavior is tested extensively in \texttt{patch class} cases in \texttt{src/\_\_tests\_\_/patch-dom.test.ts}.

\subsection{Styles: patchStyles uses objectsDiff}

\texttt{patchStyles} again uses \texttt{objectsDiff} and applies:

\begin{itemize}
  \item remove: \texttt{removeStyle(el, key)}
  \item add/update: \texttt{setStyle(el, key, value)}
\end{itemize}

Test group: \texttt{patch style}.

\subsection{Events: patchEvents}

\texttt{patchEvents} compares old and new \texttt{on} objects using \texttt{objectsDiff} and performs two phases:

\begin{enumerate}
  \item \textbf{Detach} removed and updated handlers by removing their old wrapper listeners.
  \item \textbf{Attach} added and updated handlers using \texttt{addEventListener} from \texttt{src/events.ts}.
\end{enumerate}

Why detach+attach for \emph{updated}?
Because the wrapper listener function identity changes when the handler changes.

\paragraph{Why this detail matters so much}
When people build their first patcher, they often try to "update" an event handler by doing nothing:
``the new handler is different, but the event name is the same, so surely the browser will call the new one.''

That doesn't work.
In the DOM, \texttt{addEventListener} doesn't register "an event called click".
It registers a \emph{pair}: \texttt{(eventName, functionObject)}.
Removal uses the same pair.

Relax solves this with the wrapper identity trick from Chapter~7, and the patcher follows the same rule:

\begin{itemize}
  \item if a handler is removed or updated, remove the old wrapper function
  \item if a handler is added or updated, add a fresh wrapper function
  \item store the wrapper map on the vnode (\texttt{vdom.listeners}) so destroy can remove correctly
\end{itemize}

If you look at the real implementation in \texttt{src/patch-dom.ts}, the important part is the ordering:
we remove first, then add. That prevents accidental double-binding during an update.

\paragraph{Diagram: event patching as identity management}
This is the subtle part: the browser removes listeners by the \emph{exact function object} you originally passed.

\begin{center}
\begin{tikzpicture}[
  node distance=9mm,
  box/.style={draw,rounded corners,fill=RelaxLight,inner sep=5pt,align=left},
  arrow/.style={-Latex,thick}
]
  \node[box] (old) {old props:\\\texttt{on: \{ click: oldHandler \}}};
  \node[box, right=16mm of old] (wrap1) {mount/patch created\\wrapper \texttt{w1}};
  \node[box, below=of old] (new) {new props:\\\texttt{on: \{ click: newHandler \}}};
  \node[box, right=16mm of new] (wrap2) {needs new\\wrapper \texttt{w2}};

  \node[box, below=14mm of wrap1] (detach) {detach:\\\texttt{el.removeEventListener('click', w1)}};
  \node[box, below=of wrap2] (attach) {attach:\\\texttt{el.addEventListener('click', w2)}};

  \draw[arrow] (old) -- (wrap1);
  \draw[arrow] (new) -- (wrap2);
  \draw[arrow] (wrap1) -- (detach);
  \draw[arrow] (wrap2) -- (attach);
\end{tikzpicture}
\end{center}

Tests like \texttt{patch event handlers} effectively specify this: "updating" a handler is implemented as remove+add because the wrapper identity changes.

Tests to read:

\begin{itemize}
  \item \texttt{patch event handlers} in \texttt{src/\_\_tests\_\_/patch-dom.test.ts}
  \item The binding rules in \texttt{src/\_\_tests\_\_/events.test.ts}
\end{itemize}

\subsection{Component nodes: patchComponent}

\texttt{patchComponent(oldVdom, newVdom)} preserves the component instance and pushes changes into it.

Key steps:

\begin{itemize}
  \item Extract external children and props (see \texttt{extractPropsAndEvents} in \texttt{src/utils/props.ts}).
  \item Call \texttt{component.setExternalContent(children)}.
  \item Call \texttt{component.updateProps(props)}.
  \item Copy instance fields:
    \begin{itemize}
      \item \texttt{newVdom.component = component}
      \item \texttt{newVdom.el = component.firstElement}
    \end{itemize}
\end{itemize}

This is how Relax achieves the "component instance is preserved" semantics tested in \texttt{patch component with props}.

\subsection{HRBR nodes: patchHrbr}

\subsection{A tiny but important helper: \texttt{patchDOMSameElement}}
You will notice a second patch entry point inside \texttt{src/patch-dom.ts}:

\begin{quote}
    exttt{patchDOMSameElement(oldVdom, newVdom, hostComponent?)}
\end{quote}

It is identical to \texttt{patchDOM} except it \emph{skips the replacement path} and therefore does not need \texttt{parentEl}.

This helper is used in child diffing paths where:

\begin{itemize}
  \item we already know the node should be kept (for example, keyed nodes that exist in both old and new), and
  \item we want to avoid scanning \texttt{parentEl.childNodes} to find the current index for every child.
\end{itemize}

When you're reading patching code, seeing \texttt{patchDOMSameElement} is a clue: ``we're in an optimization path, and we've already committed to reusing the node identity.''

Relax contains an experimental HRBR interop mode. In \texttt{patch-dom.ts}, \texttt{patchHrbr(old, next)}:

\begin{itemize}
  \item Reuses the existing mounted instance when \texttt{old.mount === next.mount}.
  \item Otherwise, disposes/destroys the old instance and remounts into the same host element.
\end{itemize}

This behavior matches the spec tests in \texttt{src/\_\_tests\_\_/hrbr-integration.test.ts}.

\section{Patching children: from fast paths to general diff}

After patching the node itself, Relax patches children with:

\begin{lstlisting}[language=JavaScript]
patchChildren(oldVdom, newVdom, hostComponent)
\end{lstlisting}

This is where most performance work lives, and \texttt{src/patch-dom.ts} contains several layered strategies.

\subsection{Step 0: child extraction}

Children are normalized by \texttt{extractChildren(oldVdom)} / \texttt{extractChildren(newVdom)} from \texttt{src/h.ts}.
This lets fragments and components present children uniformly.

\subsection{Fast path A (opt-in): delegate keyed list reconciliation to HRBR reconciler}

There is an optional integration with the runtime reconciler:

\begin{verbatim}
if (newVdom.props.\_reconcile === 'hrbr') { ... reconcileChildren(...) }
\end{verbatim}

This path only applies when \emph{all direct children are keyed} (keyed list). It builds \texttt{ReconcileNode} specs and lets \texttt{reconcileChildren} (from \texttt{runtime/reconciler}) patch/move nodes with a single keyed pass.

You can think of it as: "VDOM delegates list diffing to a specialized keyed reconciler".

\subsection{Fast path B (opt-in): \_textOnly rows}

For large lists of rows shaped like \texttt{<li>{label}</li>}, Relax supports an opt-in hint:

\begin{itemize}
  \item If \texttt{oldVdom.props.\_textOnly === true} and \texttt{newVdom.props.\_textOnly === true}, and
  \item both old and new children are exactly one TEXT node,
\end{itemize}

then patchChildren patches only the text node via \texttt{patchDOMSameElement}.

Test name: \texttt{patchChildren row fast path (\_textOnly)}.

\subsection{Fast path C: keyed reorder via LIS}

The LIS (Longest Increasing Subsequence) trick is a classic VDOM strategy:

\begin{itemize}
  \item If you can identify the subsequence of children that are already in correct relative order,
  \item then you only need to move the remaining nodes.
\end{itemize}

Relax's implementation is in \texttt{patchKeyedChildrenReorder(...)} and uses an internal helper \texttt{lisIndices(arr)} which computes LIS indices over a sequence of old positions (ignoring \texttt{-1} markers for new nodes).

One very practical detail from the code: \texttt{elementsForChild(child)} returns \emph{multiple} DOM elements for components (using \texttt{child.component.elements}) but only one for normal elements. That's how the reorder path correctly moves a component that renders multiple siblings.

\texttt{patchKeyedChildrenReorder(oldChildren, newChildren, parentEl, hostComponent)} implements keyed reordering with minimal DOM moves:

\begin{enumerate}
  \item Ensure both sides are fully keyed.
  \item Remove old nodes not present in the new key set.
  \item Patch existing nodes and mark new nodes with \texttt{-1} in a sequence.
  \item Compute the Longest Increasing Subsequence (LIS) over old indices.
  \item Walk backwards, mounting new nodes and moving only nodes not in LIS.
\end{enumerate}

This is a classic strategy: nodes in the LIS are already in relative order, so moving the rest yields the fewest inserts/reorders.

\subsubsection*{Diagram: keyed children reconciliation with LIS}
This is the bird's-eye view of what \texttt{patchKeyedChildrenReorder} is doing:

\begin{center}
\begin{tikzpicture}[
  node distance=10mm,
  box/.style={draw,rounded corners,fill=RelaxLight,inner sep=5pt,align=center},
  arrow/.style={-Latex,thick}
]
  \node[box] (old) {Old children\\keys: A\;B\;C\;D\;E};
  \node[box, right=18mm of old] (new) {New children\\keys: B\;E\;A\;D};

  \node[box, below=of old] (map) {Map new keys\\to old indices\\(missing \textrightarrow{} -1)};
  \node[box, below=of map] (seq) {Sequence of old indices\\for new order\\e.g. [1,4,0,3]};
  \node[box, below=of seq] (lis) {Compute LIS indices\\(nodes already in\\relative order)};
  \node[box, below=of lis] (walk) {Walk backwards\\Mount new nodes\\Move non-LIS nodes};

  \draw[arrow] (old) -- (map);
  \draw[arrow] (new) -- (map);
  \draw[arrow] (map) -- (seq);
  \draw[arrow] (seq) -- (lis);
  \draw[arrow] (lis) -- (walk);
\end{tikzpicture}
\end{center}

The reason for the backward walk is practical DOM surgery: when you insert/move nodes from the end, you don't invalidate the insertion positions you haven't processed yet.

Test names to read:

\begin{itemize}
  \item \texttt{patchChildren keyed reorder fast path (LIS)}
  \item The many keyed reorder cases in \texttt{src/\_\_tests\_\_/patch-dom.test.ts}
\end{itemize}

\subsection{Fast path D: stable keyed same-order}

If keys are present and the order is identical, Relax avoids diff work entirely:

\begin{verbatim}
if (isStableKeyedSameOrder(oldChildren, newChildren)) {
  for (i) patchDOMSameElement(oldChildren[i], newChildren[i])
  return
}
\end{verbatim}

Test name: \texttt{patchChildren fast path (stable keyed order)}.

\subsection{General path: arraysDiffSequence}

If none of the fast paths apply, Relax computes an edit script:

\begin{verbatim}
const diffSeq = arraysDiffSequence(oldChildren, newChildren, areNodesEqual)
\end{verbatim}



	exttt{arraysDiffSequence} is a general-purpose utility in \texttt{src/utils/arrays.ts}. It's worth reading because it is one of the algorithmic cores of Relax.

It has extensive unit tests in:

\begin{itemize}
  \item \texttt{src/\_\_tests\_\_/diff-sequence.test.ts}
  \item \texttt{src/\_\_tests\_\_/diff-sequence.fuzzy.test.ts}
\end{itemize}

It produces operations of types \texttt{ADD}, \texttt{REMOVE}, \texttt{MOVE}, \texttt{NOOP} (see \texttt{ARRAY\_DIFF\_OP}).

\subsubsection*{How to read the diff algorithm (from scratch)}


	exttt{arraysDiffSequence(oldArr, newArr, equals)} returns an \emph{edit script}.
An edit script is just a list of operations that, when applied left-to-right, transforms \texttt{oldArr} into \texttt{newArr}.

Relax's implementation in \texttt{src/utils/arrays.ts} is practical, not academic:

\begin{itemize}
  \item It walks \texttt{newArr} from left to right.
  \item At each index it decides whether the current old item must be removed, can stay (NOOP), needs to be added (ADD), or needs to be moved (MOVE).
  \item It tracks \emph{original indices} so tests can assert stable semantics even as the working array mutates.
\end{itemize}

This is exactly the kind of algorithm that's easy to get "mostly working" and still be wrong in edge cases.
That's why Relax has both:

\begin{itemize}
  \item exact expected sequences in \texttt{diff-sequence.test.ts}, and
  \item randomized fuzz tests in \texttt{diff-sequence.fuzzy.test.ts}.
\end{itemize}

I recommend reading one fuzzy test before you touch the diff logic.
It teaches you what you're \emph{actually} allowed to change.

Then \texttt{patchChildren} executes them:

\begin{itemize}
  \item \textbf{ADD}: \texttt{mountDOM(item, parentEl, index + offset, hostComponent)}
  \item \textbf{REMOVE}: \texttt{destroyDOM(item)}
  \item \textbf{MOVE}: move DOM nodes (component = multiple elements) and patch the moved subtree
  \item \textbf{NOOP}: patch the subtree in place
\end{itemize}

\subsubsection*{How Relax executes the edit script (the part people usually hand-wave)}

This is the section that makes the whole chapter ``click''.

At this point we have:

\begin{itemize}
  \item a concrete \texttt{parentEl} (the real DOM element we will mutate),
  \item \texttt{oldChildren} (already mounted; each node has \texttt{.el} pointers),
  \item \texttt{newChildren} (fresh objects from render; \texttt{.el} mostly unset), and
  \item \texttt{diffSeq}, a list of ops that converts old \textrightarrow{} new.
\end{itemize}

The important subtlety is this: \textbf{the DOM has one truthy linear order, but our algorithm temporarily has two competing linear orders}.
As we apply operations, \texttt{oldChildren} is being mutated (conceptually) and the DOM is being mutated (physically).
That is why \texttt{patchChildren} needs an \texttt{offset}.

\paragraph{What is \texttt{offset}?}

Think of \texttt{index} in the edit script as an index in the \emph{new} children list.
But when we insert or remove DOM nodes, the ``current DOM index we should insert at'' shifts relative to that new index.

Relax models that shift with:

\begin{itemize}
  \item \texttt{offset += 1} when it mounts an item (the DOM got longer than the old list),
  \item \texttt{offset -= 1} when it destroys an item (the DOM got shorter than the old list).
\end{itemize}

So the effective DOM insertion index becomes \texttt{index + offset}.

\subsubsection*{Operation semantics in detail}

Below is a readable version of what the real code is doing (names simplified; still faithful to intent):

\begin{lstlisting}[language=JavaScript]
let offset = 0

for (const op of diffSeq) {
  const { type, index, item, from } = op

  if (type === 'ADD') {
    // Insert a brand-new subtree.
    mountDOM(item, parentEl, index + offset, hostComponent)
    offset += 1
    continue
  }

  if (type === 'REMOVE') {
    // Tear down a mounted subtree (also unbinds events).
    destroyDOM(item)
    offset -= 1
    continue
  }

  if (type === 'MOVE') {
    // 1) Move DOM nodes to the new position
    // 2) Then patch the subtree in place (so props/events update)
    // (MOVE is not just reorder; it implies we still need to patch.)
    const oldChild = oldChildren[from]
    const newChild = item

    moveDomForVNode(oldChild, parentEl, index + offset)
    patchDOMSameElement(oldChild, newChild, hostComponent)
    continue
  }

  // NOOP: same identity at this position, but props may differ.
  patchDOMSameElement(item.old, item.new, hostComponent)
}
\end{lstlisting}

This is also why Relax is strict about \texttt{areNodesEqual}: if a node is ``equal'', we assume we can move/patch it without changing its fundamental DOM shape.

\paragraph{MOVE for components (multi-element reality)}

Moving an \emph{element vnode} means moving one DOM element.
Moving a \emph{component vnode} often means moving \textbf{multiple} DOM siblings, because a component can render "fragment-like" output.

That is why the implementation builds the move list like:

\begin{lstlisting}[language=JavaScript]
const elementsToMove = isComponent(oldChild)
  ? oldChild.component.elements
  : [oldChild.el]
\end{lstlisting}

Then it inserts each element in order at the target position.
If you only moved \texttt{component.firstElement} and forgot the rest, you'd silently corrupt the DOM order.

\subsubsection*{A concrete mental simulation}

Here is a small sequence you can replay by hand:

\begin{itemize}
  \item old keys: \texttt{[A, B, C]}
  \item new keys: \texttt{[B, D, A]}
\end{itemize}

A sensible edit script is:

\begin{itemize}
  \item MOVE \texttt{B} to index 0
  \item ADD \texttt{D} at index 1
  \item MOVE \texttt{A} to index 2
  \item REMOVE \texttt{C}
\end{itemize}

As soon as \texttt{D} is added, the DOM becomes longer than the original, so the correct target indices for later operations must account for that. That's exactly what \texttt{offset} encodes.

\subsection{What the patch-children tests specify (without running them)}

The tests in \texttt{src/\_\_tests\_\_/patch-dom.test.ts} are effectively a spec for these semantics. In particular, they lock down:

\begin{itemize}
  \item \textbf{No-op patching still updates}: even when the vnode identity matches, prop diffs still apply.
  \item \textbf{Moves preserve identity but update content}: moved nodes keep their DOM element(s), but still get patched.
  \item \textbf{Component move moves all owned elements}: multi-element components reorder correctly.
  \item \textbf{Index stability}: when replacing the root, new nodes mount into the same parent index.
\end{itemize}

If you ever change \texttt{patchChildren}, read those tests first. They tell you what behavior is intentional and what behavior would be a breaking change.

\paragraph{Component moves}

A subtle point: moving a component means moving \emph{all} DOM elements owned by that component.
That’s why the code does:

\begin{verbatim}
const elementsToMove = isComponent(oldChild)
  ? oldChild.component.elements
  : [oldChild.el]
\end{verbatim}

This matches the ``component as a fragment'' mental model: the component is a logical node that may span multiple DOM siblings.

\subsection{What the diff-sequence tests prove}
The patcher relies on \texttt{arraysDiffSequence} to produce a valid edit script. The contract is captured in two test files:

\begin{itemize}
  \item \texttt{src/\_\_tests\_\_/diff-sequence.test.ts}: a human-readable set of exact expected operation sequences for many small cases
  \item \texttt{src/\_\_tests\_\_/diff-sequence.fuzzy.test.ts}: a property-style test that generates random arrays and asserts that applying the operations reproduces the target array
\end{itemize}

The workhorse assertion is essentially:

\begin{quote}
    exttt{applyArraysDiffSequence(oldArray, arraysDiffSequence(oldArray, newArray)) === newArray}
\end{quote}

This is a great example of how Relax uses tests: not just to assert a single result, but to lock down the correctness of an algorithm under many variations.

\section{Tests to lean on}

If you want to understand Relax patch semantics quickly, scan these groups in \texttt{src/\_\_tests\_\_/patch-dom.test.ts}:

\begin{itemize}
  \item Root replacement and \texttt{.el} preservation.
  \item Fragment patching (nested add/remove/move).
  \item Attribute/class/style patching.
  \item Event handler updates and removals.
  \item Children patching: keyed/unkeyed, moves, fast paths.
  \item Component patching and component list semantics.
\end{itemize}

\section{Takeaways}

\begin{itemize}
  \item Relax’s patcher is layered: correct general diff + several pragmatic fast paths for real-world performance.
  \item Keys are not just an optimization; they unlock better move semantics and fewer DOM operations.
  \item Correctness depends on two invariants:
    \begin{itemize}
      \item \texttt{areNodesEqual} must be consistent across mount/patch.
      \item \texttt{newVdom.el} (and component element arrays) must be updated to reflect reality.
    \end{itemize}
\end{itemize}

\section{Online resources}
\begin{itemize}
  \item React keys and list rendering (conceptual background; implementation differs): \href{https://react.dev/learn/rendering-lists}{https://react.dev/learn/rendering-lists}
  \item Longest Increasing Subsequence (LIS) algorithm overview: \href{https://en.wikipedia.org/wiki/Longest\_increasing\_subsequence}{https://en.wikipedia.org/wiki/Longest\_increasing\_subsequence}
  \item DOM operations cost model (MDN on \texttt{Node.insertBefore} / \texttt{Node.removeChild}): \href{https://developer.mozilla.org/en-US/docs/Web/API/Node/insertBefore}{https://developer.mozilla.org/en-US/docs/Web/API/Node/insertBefore}
  \item DOM tree mutation semantics (MDN on \texttt{Node.replaceChild}): \href{https://developer.mozilla.org/en-US/docs/Web/API/Node/replaceChild}{https://developer.mozilla.org/en-US/docs/Web/API/Node/replaceChild}
  \item Virtual DOM diffing background reading (generic): \href{https://increment.com/frontend/a-deep-dive-into-virtual-dom/}{https://increment.com/frontend/a-deep-dive-into-virtual-dom/}
\end{itemize}
